% ============================================================
% discussion.tex — General Discussion
% ============================================================

\section{General Discussion}
\label{sec:discussion}

This study examined cortical morphometry across two large, well-phenotyped cohorts with the goal of characterizing how sleep quality, fluid intelligence, neuroticism, and age relate to cortical thickness, surface area, and volume.
The analyses produced three principal contributions, which we discuss in turn below.

\subsection{A Developmental Dissociation for Neuroticism}

Perhaps the most striking finding is the divergence between the young-adult and aging cohorts in the structural correlates of neuroticism.
In the HCP sample (ages 22--37), neuroticism was associated with reduced cortical surface area (39 of 68 DK regions) but was entirely unrelated to cortical thickness (0 of 68 regions).
In the AABC sample (ages 36--89), neuroticism was associated with reduced cortical thickness in 327 of 360 Glasser regions, with the largest effects in the anterior and posterior cingulate cortex, medial prefrontal cortex, and the parieto-occipital sulcus.

This pattern suggests that the neuroanatomical signature of neuroticism has two components: a developmentally stable component reflected in surface area, which is largely established in early life and shows high heritability~\cite{grasby2020genetic}, and a time-dependent component reflected in thickness, which emerges over decades and may reflect the cumulative effect of chronic stress on cortical tissue~\cite{lupien2009effects, mcewen2006protective}.
The biological pathway most likely to mediate the latter component involves the HPA axis.
Individuals high in neuroticism tend to produce more cortisol in response to daily stressors, exhibit flatter diurnal cortisol slopes, and show elevated inflammatory markers~\cite{ottaviani2016physiological}.
Over years, this chronic physiological activation could produce dendritic retraction, synaptic loss, and ultimately a measurable reduction in cortical thickness~\cite{lupien2009effects}, particularly in regions rich in glucocorticoid receptors such as the cingulate and medial prefrontal cortex~\cite{mcewen2006protective}.

This dissociation has methodological implications.
Studies that examine personality and brain structure exclusively in young adults may conclude that neuroticism has no thickness-related structural correlate, which is technically accurate but potentially misleading.
Our results suggest that the absence of a thickness effect in young adults does not mean that the biological processes linking neuroticism to cortical structure are absent; rather, they may not yet have accumulated to a detectable level.

\subsection{The Dominance of Age in Cortical Morphometry}

The aging analyses revealed that age was negatively correlated with cortical thickness in 356 of 360 Glasser regions and with cortical volume in 355 of 360 regions.
The single largest correlation (\textit{r}\,=\,$-$.69 for primary motor cortex thickness) implies that age explains approximately 48\% of the between-subject variance in that region.
This degree of explanatory power is exceptional for a single predictor in a cross-sectional neuroimaging study and highlights the dominance of age as a determinant of cortical structure.

The practical consequence of this dominance is that behavioral predictors with modest effect sizes may be undetectable in age-heterogeneous samples.
This is precisely what we observed for sleep quality: a variable that produced Cohen's \textit{d} values of 0.20--0.25 in the age-restricted HCP sample showed no surviving associations in the AABC sample, where the age-driven variance eclipsed the sleep-related signal.
Researchers who wish to study brain-behavior associations in aging populations would benefit from either statistically partitioning the age signal (e.g., via partial correlation or regression residualization) or focusing on longitudinal change, where within-person variation can be separated from between-person age differences~\cite{fjell2023poor, storsve2014differential}.

\subsection{Fluid Intelligence and Cortical Surface Area}

The universal positive association between fluid intelligence and cortical surface area across all 68 DK regions reinforces the view that intelligence is a brain-wide property rather than a localized function~\cite{deary2010neuroscience}.
The finding that the strongest correlations were in temporal association cortex, rather than in the parietal regions emphasized by the P-FIT model~\cite{jung2007pfit}, adds nuance to existing theory.
The inferior and middle temporal gyri support visual object recognition and semantic integration, processes that are directly engaged by the matrix-reasoning task used to measure fluid intelligence in the HCP battery.
It is possible that the temporal predominance reflects task-specific demands rather than a general property of the intelligence--surface area relationship.
Studies using more diverse cognitive batteries would be needed to test this possibility.

The specificity of the intelligence effect to surface area (rather than thickness) is noteworthy.
Surface area reflects the tangential extent of the cortical sheet and is determined primarily by the proliferation of radial progenitor cells during prenatal development~\cite{grasby2020genetic}.
Cortical thickness, by contrast, reflects the radial organization of the cortical column and is more sensitive to postnatal environmental influences.
The fact that intelligence related to surface area but not thickness (in this sample, at least) suggests that the structural foundation of fluid intelligence is rooted primarily in prenatal cortical development rather than in experience-dependent plasticity.

\subsection{Limitations}

Several limitations should be acknowledged.
First, both analyses are cross-sectional, which precludes causal inference.
The associations reported here may reflect shared genetic variance, confounding environmental factors, or reverse causation.
Second, the two cohorts use different atlases (Desikan-Killiany vs.\ Glasser HCP-MMP1), which differ in spatial resolution by a factor of approximately five.
Although we constructed a lobe-level correspondence table to facilitate qualitative comparison, region-level results cannot be directly compared across datasets.
Third, sleep quality and neuroticism were assessed using self-report instruments, which are subject to recall bias and may not capture the physiological dimensions of sleep or stress reactivity that are most relevant to brain structure.
Objective sleep measures (e.g., actigraphy, polysomnography) and physiological stress measures (e.g., salivary cortisol) would provide stronger tests of the proposed mechanisms.
Fourth, the HCP sample's restricted age range (22--37) is a methodological strength for isolating behavioral effects from age confounds, but it limits generalizability to the broader population.

\subsection{Future Directions}

Two directions for future research emerge from this work.
First, longitudinal analyses within both cohorts, as data become available, would allow us to test whether within-person changes in cortical thickness track changes in sleep quality or stress-related measures, providing a stronger basis for causal inference.
Second, mediation analyses incorporating salivary cortisol, inflammatory markers, or composite allostatic load indices could test the hypothesized HPA-axis pathway linking neuroticism to cortical thinning in older adults.
